\section{APRESENTAÇÃO E ANÁLISE DOS RESULTADOS}
\label{sec:apresentacao_analise_resultados}
Este capítulo detalha os procedimentos usados para testar, validar e avaliar a
qualidade do processo desenvolvido, contemplando desde os pontos fortes da
abordagem, erros encontrados e o que pode ser melhorado.

\subsection{Ambiente de execução e hardware utilizado}
\label{sec:ambiente_execucao}
Para a avaliação de desempenho, bem como a execução do \textit{pipeline} de
processamento proposto neste trabalho, foi utilizado um ambiente virtualizado
disponibilizado pela SEFAZ. A documentação desse ambiente tem por objetivo
garantir reprodutibilidade do que será apresentado e fornecer uma base para
análises futuras, especialmente no que tange ao paralelismo central da
aplicação.

Todos os processos do sistema foram executados em uma única máquina virtual
hospedada sob um \textit{hypervisor} \textit{VMWare} com virtualização total.
As especificações dos recursos alocados para esta instância encontram-se no
Quadro~\ref{tab:specs_ambiente}.

\begin{quadro}[hbt!]
    \centering
    \caption{Especificações do ambiente computacional}
    \label{tab:specs_ambiente}
    \begin{tabular}{|p{3.5cm}|p{8.5cm}|}
        \toprule
        \textbf{Componente} & \textbf{Especificação Técnica} \\
        \midrule
        \textbf{Infraestrutura} & Máquina Virtual sob \textit{hypervisor} VMware (\textit{virtualização total}) \\
        \textbf{Processador} & Intel\textsuperscript{\tiny\textregistered} Xeon\textsuperscript{\tiny\textregistered} Gold 5320 @ 2.20GHz \\
        \textbf{Topologia da CPU} & 16 núcleos lógicos (16 \textit{sockets}, 1 \textit{core/socket}, 1 \textit{thread/core}) \\
        \textbf{Cache} & L1: 80 KB / L2: 1,25 MB / L3: 39,9 MB (Compartilhado) \\
        \textbf{Memória RAM} & 93Gi, com 4Gi de swap \\
        \textbf{Armazenamento (Raiz do sistema) } & SSD de 135GB \\
        \textbf{Armazenamento (Volume dedicado ao banco de dados) } & SSD de 36TB \\
        \textbf{Versão do \textit{PostgreSQL}} & 16 \\
        \bottomrule
    \end{tabular}

    \vspace{0.2cm}
    \small{Fonte: Elaborado pelo autor (2026).}
\end{quadro}


A infraestrutura provisionada apresenta características altamente favoráveis
para o perfil de carga de trabalho do sistema desenvolvido. Destaca-se
principalmente a topologia de processamento disponibilizada: a alocação de 16
processadores lógicos concentrados em um único nó NUMA, garantindo um tempo de
acesso uniforme à memória RAM. Na prática, isso permitiu com que o Orquestrador
instanciasse até 16 \textit{Workers}, saturando a capacidade
de processamento paralelo, sem sofrer degradação devido a penalidades de
troca de contexto. O banco de dados utilizado para testes foi instanciado
na mesma máquina virtual, o que manteve os testes isolados de
eventuais flutuações relacionadas à latência de rede. 

\subsection{Metodologia de testes}
\label{sec:testes}
A validação do \textit{pipeline} de ingestão desenvolvido foi conduzido por
meio testes experimentais de caráter preliminar, em grande parte de natureza
\textit{ad-hoc}. Devido a restrições de tempo, experiência da equipe e escopo
do projeto, não foi possível a realização de uma bateria de testes ou de uma
coleta de dados em situação de estresse em ambiente de produção real.

Contudo, cumpre destacar que os cenários selecionados para os testes
\textit{ad-hoc} foram projetados estrategicamente para cobrir fluxos críticos
de execução do sistema, servindo como uma validação da arquitetura
proposta. Para tanto, a massa de dados de teste consistiu em 100 arquivos
EFD ICMS IPI reais, número este que foi escolhido pela equipe por ser
considerado um número representativo, mas não houve nenhuma avaliação da
eficácia real dessa escolha. Também foram usados diversos arquivos
EFD ICMS IPI gerados artificialmente, buscando cobrir as seguintes
especificações:
\begin{itemize}
    \item Arquivos contendo uma unidade de cada um dos registros mapeados;
    \item Arquivos com ausência de registros que foram normalizados;
    \item Arquivos cobrindo os diversos casos de um-para-vários;
    \item Arquivos com registros contendo menos campos que o total mapeado;
    \item Arquivos contendo casos de um-para-vários onde se espera um-para-um;
    \item Arquivos sem a presença do registro 0000;
    \item Arquivos com campos que quebram as regras de formatação especificada;
    \item Arquivos com campos que não quebram as regras de formatação, mas não
    conseguem ser inseridos por limitações técnicas. Ex: Campos numéricos que
    causam \textit{overflow};
    \item Arquivos com registros contendo mais campos que o total esperado;
\end{itemize}
Esta composição visou cobrir possíveis lacunas e anomalias que
os documentos reais pudessem deixar de apresentar.

Uma das estratégias adotadas para maximizar a eficácia dos testes
\textit{ad-hoc} concentrou-se na verificação da integridade absoluta da
informação trafegada ao longo do \textit{pipeline}. Empregou-se uma idéia
semelhante a técnica de \textit{round-trip validation}, que consiste na
validação dos dados mediante a reconstrução reversa da informação original a
partir dos registros já processados. Apesar da equipe não construir um
processo totalmente automatizado e rico como o esperado pelo método
\textit{round-trip validation}, ainda foi possível verificar a premissa central
de que, se o sistema é capaz de desestruturar, transformar e persistir os dados
corretamente, deve ser igualmente capaz de remontá-los em seu formato bruto
original, sem qualquer depreciação. O procedimento foi dividido em duas etapas
de verificação isoladas, valendo-se de análises diferenciais de texto:

\begin{enumerate}
    \item \textbf{Validação em Memória (\textit{Parser}):} O arquivo
    EFD ICMS IPI original foi processado em memória e convertido para o formato
    estruturado em JSON. Em seguida, aplicou-se uma rotina de engenharia
    reversa para reconverter esse JSON de volta para o formato de texto plano
    original. Além disso, o arquivo EFD ICMS IPI original foi normalizado para
    evitar divergências devido à representações equivalentes, como costuma
    ocorrer entre campos numéricos. A comparação entre o arquivo de origem e o
    gerado a partir do JSON atestou a precisão da modelagem estrutural
    em árvore;
    \item \textbf{Validação de Persistência (\textit{Loader}):} Para homologar
    a carga no banco de dados, o ciclo de ingestão foi concluído com a inserção
    dos registros no \textit{PostgreSQL}. Posteriormente, utilizando
    exclusivamente as chaves primárias para rastrear os relacionamentos
    hierárquicos entre as tabelas, os registros foram consultados, extraídos
    e remontados de volta, no formato original, seguindo o mesmo processo do
    caso anterior a partir daqui.
\end{enumerate}

Em ambos os cenários de teste, a verificação diferencial linha a linha confirmou
a exatidão do processo. A ausência de discrepâncias relevantes provou
empiricamente que o sistema cumpre o requisito de extração e transformação,
sem incorrer em perdas de informação, truncamento de valores, ou corrupção na
hierarquia de dependência dos registros fiscais.

Outra abordagem consistiu na elaboração de documentos contendo as falhas de
construção pré-determinadas listadas anteriormente, inserindo-os nos lotes de
processamento para avaliação de comportamento. Este ensaio buscou aferir a
resiliência do sistema no tratamento de arquivos defeituosos e avaliar a
capacidade de rastreabilidade dos erros, permitindo diagnosticar se a falha
advém do arquivo, de equívocos nos metadados ou da lógica de código.

Neste cenário, os resultados mostraram-se satisfatórios. Os \textit{logs}
gerados apresentaram um nível adequado de detalhamento
(Listagem~\ref{qua:log_parsing_incorreto}). A identificação explícita do erro,
associada ao nome do arquivo, viabiliza a averiguação direta da anomalia.
Adicionalmente, constatou-se que a ocorrência não comprometeu o
\textit{pipeline}, o qual prosseguiu sua execução, sem sofrer interrupções
globais. Este comportamento atesta empiricamente que o sistema isola arquivos
defeituosos, otimizando a utilização dos recursos alocados para o
respectivo lote.

\begin{quadro}[H]
    \centering
    \caption{Exemplo de log de erro de \textit{parsing}}
    \label{qua:log_parsing_incorreto}
\begin{lstlisting}[basicstyle=\ttfamily\footnotesize, frame=single]
12:11:04,445 INFO ERROR processing EFDs/F0000005_AAA+AAAA_EFD_999-9999-MOCK-99999999999999-999999999-0-202512-17122025144834-999-17122025144834.txt.zip: Error: 0005 marked as single occurrence per parent, but more than one appeared:
	Register:|0005|EMPRESA TESTE LTDA|EMPRESA |COMERCIAL ABC|EMPRESA TE|FULANO DE TAL|RUA DAS FLORES 123|EMPRESA TES|COMERCIAL A|COMERCIAL ABC|
\end{lstlisting}

    \vspace{0.2cm}
    \small{Fonte: Elaborado pelo autor (2026).}
\end{quadro}


A estratégia final compreendeu a injeção de erros artificiais durante a execução
do \textit{pipeline}, especificamente durante a execução do componente
\textit{Loader}. Os erros considerados foram:
\begin{itemize}
    \item Falha de conexão com o \textit{PostgreSQL};
    \item Inconformidade entre o dado a ser inserido e o tipo esperado;
    \item Datas inválidas.
\end{itemize}
Tais erros foram atrelados apenas ao processamento de determinados arquivos,
com o objetivo de validar a atomicidade da carga diante de falhas no meio do
processo e verificar a precisão do registro do erro.

Neste caso, os resultados também foram positivos. Conforme os \textit{logs}
(Listing~\ref{qua:log_loading_incorreto}), a nomeação padronizada da falha e do
arquivo proporcionou boa rastreabilidade. Validou-se que todos os demais
arquivos do lote atravessaram o fluxo de carga com sucesso, enquanto o arquivo
defeituoso não produziu qualquer efeito no banco de dados. Isso comprova que o
sistema mantém a atomicidade das operações, revertendo inserções parciais
por meio de \textit{rollback} quando anomalias interrompem o fluxo.

\begin{quadro}[H]
    \centering
    \caption{Exemplo de log de erro durante a carga}
    \label{qua:log_loading_incorreto}
\begin{lstlisting}[basicstyle=\ttfamily\footnotesize, frame=single]
DETAIL:  A field with precision 15, scale 2 must round to an absolute value less than 10^13.
13:30:53,49 INFO ERROR processing EFDs/F0000049_AAA+AAAA_EFD_999-9999-MOCK-99999999999999-999999999-0-202512-17122025144834-999-17122025144834.txt.zip: DB insert error for b020: numeric field overflow
\end{lstlisting}

    \vspace{0.2cm}
    \small{Fonte: Elaborado pelo autor (2026).}
\end{quadro}

Portanto, conclui-se que a avaliação do sistema atendeu com êxito a três
critérios fundamentais:
\begin{enumerate}
    \item \textbf{Consistência do componente \textit{Parser}:} Verificação de
    que a árvore de registros gerada em memória refletiu com fidelidade a
    hierarquia descrita na especificação fiscal, abrangendo as particularidades
    estruturais e processando todos os dados;
    \item \textbf{Aderência aos Metadados:} Validação da capacidade de
    interpretar dinamicamente o arquivo JSON de regras, assegurando sua
    aplicação, isolando arquivos não conformes e alternando entre os modos
    de inserção adequados;
    \item \textbf{Atomicidade da Carga:} Confirmação de que o componente
    \textit{Loader} realizou a persistência normalizada de maneira íntegra,
    respeitando as restrições de integridade referencial e executando
    \textit{rollback} imediato em caso de falhas, protegendo o SGBD contra
    dados corrompidos.
\end{enumerate}

Embora os dados coletados não permitam traçar curvas generalizáveis de
escalabilidade linear sob concorrência extrema de \textit{Workers}, os
resultados obtidos atestam a viabilidade técnica da solução e a corretude
lógica de todos os módulos da arquitetura.

\subsection{Limitações encontradas após a modelagem}
\label{sec:limitacoes}

As subseções a seguir detalham gargalos e desafios arquiteturais identificados
após a modelagem e durante as etapas de execução do \textit{pipeline}.

\subsubsection{Limitações do Tipo JSONB com Grandes Volumes}
\label{sec:json_grande}
Após a conclusão dos ensaios preliminares, o sistema foi avaliado em um
ambiente de homologação controlado, visando observar seu comportamento
perante um volume de dados maior e mais condizente com cenários reais.
Um dos imprevistos ocorridos esteve diretamente relacionado às restrições
físicas de armazenamento do SGBD \textit{PostgreSQL} no tratamento de
arquivos EFD ICMS IPI de grande escala. 

O projeto original previa que, em paralelo à carga relacional normalizada, o
\textit{Worker} realizaria a conversão do arquivo completo para o
formato JSON bruto (\textit{raw data}) para persistência em uma única
coluna do tipo \texttt{JSONB}. Contudo, ao submeter arquivos de
volumetria massiva (centenas de milhares de registros), o SGBD disparou
erros de estouro de limite físico, interrompendo as transações. Esses
arquivos chegaram a ocasionar falhas no próprio \textit{pool}\footnote{Conjunto
de conexões ativas e reutilizáveis mantido em memória pelo sistema, evitando
o custo computacional de abrir e fechar novas conexões com o banco a cada
operação.} de conexões com o banco, resultando no cancelamento da inserção de
lotes inteiros, mesmo com a configuração de \textit{retries}\footnote{Estratégia
de resiliência que agenda a reexecução automática da criação de uma conexão com
o banco de dados antes de declará-lo definitivamente fora de alcance.}
habilitada.

Essa falha decorre dos limites estáticos impostos pelo \textit{PostgreSQL}
para a alocação de objetos JSONB. O erro mais recorrente foi a violação
do limite de 256 MB para campos únicos, embora também tenham sido registradas
instâncias de estouro do limite máximo de 1 GB estipulado pelo SGBD.

Para mitigar essa limitação, determinou-se que a inserção de arquivos brutos
monolíticos deve ser evitada em casos de alta volumetria. Incorporou-se uma
cláusula de verificação que estima o tamanho do JSON gerado antes da operação
de E/S --- Sigla para \textit{Entrada e Saída}, que representa operações que
acessam recursos externos ao programa, como o SGBD. Caso a estimativa ultrapasse
o limite configurado, a persistência no campo JSONB é suprimida,
garantindo que o processamento normalizado avance sem interrupções. 

A descoberta deste gargalo evidenciou que a persistência de superobjetos
estruturados em SGBDs relacionais representa um anti-padrão de engenharia
para lidar com volumes característicos de \textit{Big Data}\footnote{Termo que
descreve conjuntos de dados massivos, complexos e de rápido crescimento, cuja
escala e variedade ultrapassam a capacidade de armazenamento, gestão e
processamento de SGBDs relacionais convencionais.}, gerando
conhecimento valioso para a evolução arquitetural da ferramenta e
do modelo de dados utilizado pela equipe.

\subsubsection{Degradação da inserção}
\label{sec:degradacao_insercao}
Durante a validação do componente \textit{Loader} com arquivos EFD ICMS IPI
massivos, observou-se uma degradação não linear no tempo total de persistência.
Embora a ingestão de dados em tempo real não figure entre
os requisitos não funcionais deste trabalho, a análise dessa latência apontou
para um gargalo intrínseco à sistemas de ingestão de dados massivos.

A primeira hipótese para a raiz do comportamento foi a dependência estrutural
de chaves primárias artificiais geradas sequencialmente pelo
\textit{PostgreSQL}. Devido à natureza hierárquica da obrigação fiscal, a
persistência exige a inserção prévia do registro-pai para obtenção de seu
identificador único gerado pelo banco. Somente após receber este valor, o
componente \textit{Loader} é capaz de propagá-lo como chave estrangeira
aos registros-filhos.

Contudo, as investigações demonstraram que a degradação de desempenho observada
não decorre unicamente dessa dependência. Uma remodelagem experimental
eliminando a necessidade de consultas prévias de identificadores foi
implementada e testada, porém o comportamento de perda progressiva de vazão
persistiu de forma análoga.

A análise empírica refinada nos levou a concluir que a verdadeira causa
raiz reside na sobrecarga imposta ao motor de persistência do
\textit{PostgreSQL}, durante a inserção de volumes massivos sob um único bloco
transacional. À medida que milhões de registros hierárquicos são injetados
continuamente, o banco de dados é forçado a manter em tempo real a consistência
das restrições referenciais e a integridade de múltiplos índices estruturados.
Esse processo de atualização dinâmica satura os recursos de E/S e degrada o
tempo de execução de forma não-linear, conforme o volume do documento cresce.

A desativação momentânea das restrições de consistência para
a realização da carga atestou o ganho de desempenho, diminuindo
consideravelmente o tempo necessário para a execução total da carga. Contudo,
a resolução definitiva deste gargalo não pode contemplar apenas essa escolha,
e exige estratégia de ingestão de
alta complexidade, envolvendo uma modelagem dedicada do funcionamento do
\textit{PostgreSQL} para ingestão massiva de dados em único bloco transacional.
Diante dessa complexidade, e do fato de que a abordagem atual atende
satisfatoriamente aos requisitos de negócio, para a grande maioria dos
cenários, a implementação de um novo modelo de persistência otimizado foi
deliberadamente isolada do escopo atual, configurando-se como uma oportunidade
promissora para trabalhos futuros.

\subsubsection{Documentação conflitante}
\label{sec:documentacao_conflitante}
Durante as homologações, observou-se uma ocorrência errônea em relação à
base normativa. O registro \texttt{H020}, hierarquicamente subordinado ao
registro \texttt{H010}, encontra-se categorizado com multiplicidade
um-para-um no Guia Prático~\cite{guia_pratico_efd} oficial da EFD. Não obstante,
uma parcela ínfima dos arquivos analisados possuía
múltiplos registros \texttt{H020} para um mesmo \texttt{H010}.

Após averiguação, constatou-se que a Nota Técnica~\cite{nota_tecnica_efd}
correspondente indicava tal relacionamento como um-para-vários. Para mitigar o
impasse normativo e garantir maior flexibilidade à extração, assumiu-se a
especificidade da Nota Técnica~\cite{nota_tecnica_efd} como a diretriz válida
para os metadados do sistema. Este evento evidencia que as especificações
oficiais emitidas pelo fisco são suscetíveis a divergências internas, exigindo
processos contínuos de governança de dados e adaptação dinâmica das regras de
negócio.

\subsection{Métricas Operacionais Preliminares e Linha de Base}
\label{sec:numeros_hell_yeah}

Embora o escopo central deste trabalho resida na validação da arquitetura
proposta quanto a consistência dos dados inseridos e do uso adequado das regras
em JSON, o monitoramento das execuções no ambiente de homologação controlado
permitiu a coleta empírica de métricas operacionais que estabelecem
uma linha de base para trabalhos futuros.
Os ensaios registraram o tempo bruto de execução do \textit{pipeline}
em relação ao volume total de arquivos submetidos, englobando as
fases de orquestração, transformação estrutural e
carga no banco de dados.

A Tabela~\ref{tab:tempo_execucao} consolida os resultados quantitativos
extraídos do ambiente de homologação, considerando um Orquestrador trabalhando
com \( 8 \) \textit{Workers} simultâneos. Esse valor foi escolhido por consumir
cerca da metade do poder de processamento do ambiente disponibilizado,
garantindo que flutuações de desempenho causadas por outras atividades sendo
executadas fosse minimizado.
A Tabela~\ref{tab:estatisticas} reúne resultados estatísticos acerca dos
\textit{Workers} executados, cada um imbuído de processar
\( 100 \) arquivos EFD ICMS. O número \( 100 \) foi escolhido ao acaso e,
apesar de constituir um vetor de exploração interessante para a métrica de
desempenho, não foram realizados testes com outros valores.
\begin{table}[hbt!]
    \centering
    \caption{Tempo bruto de execução por volumetria}
    \label{tab:tempo_execucao}
    \begin{tabular}{p{2cm} p{3cm} p{2cm} p{4cm}}
        \toprule
        \textbf{Número da instância de execução} & \textbf{Número de EFDs} & \textbf{Volume Aproximado} & \textbf{Tempo Total de Execução} \\
        \midrule
        1 & 194668 arquivos & 9.6G & 21h 36m 25s \\
        2 & 200040 arquivos & 11G  & 28h 09m 31s \\
        3 & 242653 arquivos & 11G  & 21h 31m 16s \\
        \bottomrule
    \end{tabular}

    \vspace{0.2cm}
    \small{Fonte: Elaborado pelo autor (2026).}
\end{table}


\begin{table}[hbt!]
    \centering
    \caption{Estatísticas por lote (\textit{Baseline})}
    \label{tab:estatisticas}
    \begin{tabular}{p{2cm} p{2cm} p{2cm} p{2cm} p{2cm}}
        \toprule
        \textbf{Número da instância de execução} & \textbf{\textit{Worker} mais rápido} & \textbf{\textit{Worker} mais lento} & \textbf{Tempo médio} & \textbf{Desvio padrão} \\
        \midrule
        1 & 0{,}5m & 56{,}9m & 5{,}3m & 5{,}0m\\
        2 & 0{,}5m & 76{,}8m & 6{,}7m & 6{,}0m\\
        3 & 0{,}7m & 119{,}6m & 8{,}7m & 11{,}7m\\
        \bottomrule
    \end{tabular}

    \vspace{0.2cm}
    \small{Fonte: Elaborado pelo autor (2026).}
\end{table}


Destaca-se que não foram aplicados comparativos diretos com ferramentas
comerciais de mercado. Isso se deve ao fato da literatura carecer de
referenciais abertos semelhantes ao explorado no presente trabalho.

Ademais, o monitoramento ativo dos lotes forneceu um panorama sobre a
integridade da base de dados fornecida. De um total de \( 637.361 \) arquivos
submetidos ao \textit{pipeline}, o sistema diagnosticou e isolou \( 15 \)
arquivos contendo anomalias estruturais ou normativas
(aproximadamente \(0{,}002\)\% da amostra), dentre essas nenhuma envolvendo
a especificação em JSON criada. Este índice quantifica a suscetibilidade a
erros inerente às obrigações fiscais e justifica, na prática, o esforço
arquitetural investido nos mecanismos de resiliência, isolamento de
falhas e rastreabilidade discutidos nas seções anteriores.

Junto a isso, a alta variação entre o tempo de execução de cada \textit{Worker},
 chegando a duração de 119 minutos no pior caso total, exemplifica o que foi
 exibido na Seção~\ref{sec:degradacao_insercao}, onde foi comentado como
 determinados \textit{Workers} sofrem degradação no tempo de execução devido a
 arquivos massivos. Ainda assim, é válido ressaltar que ao longo dos três lotes executados, cerca de \( 90\% \) dos \textit{Workers} (\( 88\% \) nos lotes 1
 e 2, \( 93\% \) no lote 3) permaneceu a menos de um desvio padrão em torno da
 média aritmética, \textit{i.e.}, cerca de \( 90\% \) dos \textit{Workers}
 demora não mais que 20 minutos para executar todo seu \textit{pipeline}, o
 que denota uma boa estabilidade operacional, salvo os casos onde foi
 evidenciado a degradação do tempo de execução. O \textit{boxplot} exibido na
 Figura~\ref{fig:boxplot_tempos_1} permite uma visualização melhor do
 comportamento do tempo do execução do \textit{pipeline} em relação ao
 tempo de execução de cada {Worker}, destacando em cada um deles, os
 \( 5 \) \textit{Workers} com maior tempo de execução.

\subsection{Síntese dos Resultados}
\label{sec:discussao_final}

Os ensaios experimentais conduzidos demonstraram que a adoção de uma
arquitetura orientada a metadados e estruturada em arquivos JSON é
tecnicamente viável e altamente íntegro para o processamento de obrigações
fiscais complexas, como a EFD ICMS IPI. A bateria de testes, mesmo com sua
implementação embrionária, atestou a corretude e a precisão do algoritmo de
conversão, provando a ausência de perdas informacionais durante o fluxo de
desestruturação e carga.

Simultaneamente, a implementação revelou gargalos inerentes à modelagem
relacional de grandes volumes, notadamente as limitações de paginação do tipo
\texttt{JSONB} e a latência de operações de E/S em inserções hierárquicas
massivas. Tais descobertas, no entanto, não invalidam a ferramenta,
pelo contrário, balizaram decisões arquiteturais de mitigação e evidenciaram
que o sistema lida com anomalias físicas e lógicas de forma isolada, resiliente
e rastreável. De modo geral, o \textit{pipeline} cumpriu o propósito
de transformar arquivos textuais não estruturados em uma base de dados
relacional íntegra, estabelecendo uma fundação sólida para a construção
de sistemas de auditoria fiscais.

\begin{figure}[htb]
    \centering
    \caption{Distribuição dos tempos de processamento dos \textit{Workers} em minutos em cada instância de execução.}
    \begin{tikzpicture}
        \begin{axis}[
            width=0.75\textwidth,
            height=5cm,
            xlabel style={
                at={(0.5, 1.05)},
                anchor=south
            },
            xlabel={Tempo de Execução da Instância 1 (em minutos)},
            xmajorgrids,
            ytick=\empty,
            tick label style={font=\footnotesize},
            scaled x ticks=false,
            xmin=-5, xmax=130,
            xtick={0, 20, 40, 60, 80, 100, 120}
        ]

        \addplot+[
            boxplot prepared={
                draw direction=x,
                draw position=1,
                lower whisker=0.500,
                lower quartile=2.157,
                median=3.735,
                upper quartile=6.843,
                upper whisker=13.856
            },
        ]
        coordinates{
            (1,30.198)
            (1,34.023)
            (1,43.030)
            (1,55.625)
            (1,56.909)
        };

        \end{axis}
    \end{tikzpicture}

    \centering
    \begin{tikzpicture}
        \begin{axis}[
            width=0.75\textwidth,
            height=5cm,
            xlabel style={
                at={(0.5, 1.05)},
                anchor=south
            },
            xlabel={Tempo de Execução da Instância 2 (em minutos)},
            xmajorgrids,
            ytick=\empty,
            tick label
            style={font=\footnotesize},
            scaled x ticks=false,
            xmin=-5, xmax=130,
            xtick={0, 20, 40, 60, 80, 100, 120}
        ]

        \addplot+[
            boxplot prepared={
                draw direction=x,
                draw position=1,
                lower whisker=0.484,
                lower quartile=2.824,
                median=4.700,
                upper quartile=8.778,
                upper whisker=17.648
            },
        ]
        coordinates{
            (1,42.832)
            (1,48.263)
            (1,50.877)
            (1,51.802)
            (1,76.815)
        };

        \end{axis}
    \end{tikzpicture}

    \centering
    \begin{tikzpicture}
        \begin{axis}[
            width=0.75\textwidth,
            height=5cm,
            xlabel style={
                at={(0.5, 1.05)},
                anchor=south
            },
            xlabel={Tempo de Execução da Instância 3 (em minutos)},
            xmajorgrids,
            ytick=\empty,
            tick label
            style={font=\footnotesize},
            scaled x ticks=false,
            xmin=-5, xmax=130,
            xtick={0, 20, 40, 60, 80, 100, 120}
        ]

        \addplot+[
            boxplot prepared={
                draw direction=x,
                draw position=1,
                lower whisker=0.699,
                lower quartile=3.236,
                median=5.470,
                upper quartile=10.461,
                upper whisker=21.183
            },
        ]
        coordinates{
            (1,81.825)
            (1,94.230)
            (1,108.417)
            (1,114.311)
            (1,119.596)
        };

        \end{axis}
    \end{tikzpicture}

    \small{Fonte: Elaborado pelo autor (2026).}
    \label{fig:boxplot_tempos_1}
\end{figure}

\clearpage
